% Options for packages loaded elsewhere
% Options for packages loaded elsewhere
\PassOptionsToPackage{unicode}{hyperref}
\PassOptionsToPackage{hyphens}{url}
%
\documentclass[
  ignorenonframetext,
]{beamer}
\newif\ifbibliography
\usepackage{pgfpages}
\setbeamertemplate{caption}[numbered]
\setbeamertemplate{caption label separator}{: }
\setbeamercolor{caption name}{fg=normal text.fg}
\beamertemplatenavigationsymbolsempty
% remove section numbering
\setbeamertemplate{part page}{
  \centering
  \begin{beamercolorbox}[sep=16pt,center]{part title}
    \usebeamerfont{part title}\insertpart\par
  \end{beamercolorbox}
}
\setbeamertemplate{section page}{
  \centering
  \begin{beamercolorbox}[sep=12pt,center]{section title}
    \usebeamerfont{section title}\insertsection\par
  \end{beamercolorbox}
}
\setbeamertemplate{subsection page}{
  \centering
  \begin{beamercolorbox}[sep=8pt,center]{subsection title}
    \usebeamerfont{subsection title}\insertsubsection\par
  \end{beamercolorbox}
}
% Prevent slide breaks in the middle of a paragraph
\widowpenalties 1 10000
\raggedbottom
\AtBeginPart{
  \frame{\partpage}
}
\AtBeginSection{
  \ifbibliography
  \else
    \frame{\sectionpage}
  \fi
}
\AtBeginSubsection{
  \frame{\subsectionpage}
}
\usepackage{iftex}
\ifPDFTeX
  \usepackage[T1]{fontenc}
  \usepackage[utf8]{inputenc}
  \usepackage{textcomp} % provide euro and other symbols
\else % if luatex or xetex
  \usepackage{unicode-math} % this also loads fontspec
  \defaultfontfeatures{Scale=MatchLowercase}
  \defaultfontfeatures[\rmfamily]{Ligatures=TeX,Scale=1}
\fi
\usepackage{lmodern}

\usetheme[]{moon}
\ifPDFTeX\else
  % xetex/luatex font selection
\fi
% Use upquote if available, for straight quotes in verbatim environments
\IfFileExists{upquote.sty}{\usepackage{upquote}}{}
\IfFileExists{microtype.sty}{% use microtype if available
  \usepackage[]{microtype}
  \UseMicrotypeSet[protrusion]{basicmath} % disable protrusion for tt fonts
}{}
\makeatletter
\@ifundefined{KOMAClassName}{% if non-KOMA class
  \IfFileExists{parskip.sty}{%
    \usepackage{parskip}
  }{% else
    \setlength{\parindent}{0pt}
    \setlength{\parskip}{6pt plus 2pt minus 1pt}}
}{% if KOMA class
  \KOMAoptions{parskip=half}}
\makeatother


\usepackage{longtable,booktabs,array}
\usepackage{calc} % for calculating minipage widths
\usepackage{caption}
% Make caption package work with longtable
\makeatletter
\def\fnum@table{\tablename~\thetable}
\makeatother
\usepackage{graphicx}
\makeatletter
\newsavebox\pandoc@box
\newcommand*\pandocbounded[1]{% scales image to fit in text height/width
  \sbox\pandoc@box{#1}%
  \Gscale@div\@tempa{\textheight}{\dimexpr\ht\pandoc@box+\dp\pandoc@box\relax}%
  \Gscale@div\@tempb{\linewidth}{\wd\pandoc@box}%
  \ifdim\@tempb\p@<\@tempa\p@\let\@tempa\@tempb\fi% select the smaller of both
  \ifdim\@tempa\p@<\p@\scalebox{\@tempa}{\usebox\pandoc@box}%
  \else\usebox{\pandoc@box}%
  \fi%
}
% Set default figure placement to htbp
\def\fps@figure{htbp}
\makeatother





\setlength{\emergencystretch}{3em} % prevent overfull lines

\providecommand{\tightlist}{%
  \setlength{\itemsep}{0pt}\setlength{\parskip}{0pt}}



 


\makeatletter
\@ifpackageloaded{caption}{}{\usepackage{caption}}
\AtBeginDocument{%
\ifdefined\contentsname
  \renewcommand*\contentsname{Table of contents}
\else
  \newcommand\contentsname{Table of contents}
\fi
\ifdefined\listfigurename
  \renewcommand*\listfigurename{List of Figures}
\else
  \newcommand\listfigurename{List of Figures}
\fi
\ifdefined\listtablename
  \renewcommand*\listtablename{List of Tables}
\else
  \newcommand\listtablename{List of Tables}
\fi
\ifdefined\figurename
  \renewcommand*\figurename{Figure}
\else
  \newcommand\figurename{Figure}
\fi
\ifdefined\tablename
  \renewcommand*\tablename{Table}
\else
  \newcommand\tablename{Table}
\fi
}
\@ifpackageloaded{float}{}{\usepackage{float}}
\floatstyle{ruled}
\@ifundefined{c@chapter}{\newfloat{codelisting}{h}{lop}}{\newfloat{codelisting}{h}{lop}[chapter]}
\floatname{codelisting}{Listing}
\newcommand*\listoflistings{\listof{codelisting}{List of Listings}}
\makeatother
\makeatletter
\makeatother
\makeatletter
\@ifpackageloaded{caption}{}{\usepackage{caption}}
\@ifpackageloaded{subcaption}{}{\usepackage{subcaption}}
\makeatother

\usepackage{bookmark}
\IfFileExists{xurl.sty}{\usepackage{xurl}}{} % add URL line breaks if available
\urlstyle{same}
\hypersetup{
  pdftitle={Civil Society, Race, and Local Context in the 2024 Election},
  pdfauthor={Christian Hosam and Norris Davis},
  hidelinks,
  pdfcreator={LaTeX via pandoc}}


\title{Civil Society, Race, and Local Context in the 2024 Election}
\author{Christian Hosam and Norris Davis}
\date{}

\begin{document}
\frame{\titlepage}


\begin{frame}{Motivation}
\phantomsection\label{motivation}
\begin{figure}[H]

{\centering \pandocbounded{\includegraphics[keepaspectratio]{blueroseresearch.png}}

}

\caption{Blue Rose Democratic Party 2024 Post-Mortem}

\end{figure}%
\end{frame}

\begin{frame}{Motivation}
\phantomsection\label{motivation-1}
\begin{itemize}
\tightlist
\item
  We know that in 2024, there was a major shift in terms of the support
  that Donald Trump received from communities of color.
\end{itemize}

\begin{itemize}
\tightlist
\item
  What this plot also tells us is that it was centered in certain places
  around the country.
\end{itemize}

\begin{itemize}
\tightlist
\item
  What is harder to tell is what, if anything, about those communities
  and their local contexts might have contributed to this shift.
\end{itemize}
\end{frame}

\begin{frame}{Civic Enmeshment}
\phantomsection\label{civic-enmeshment}
\begin{itemize}
\tightlist
\item
  Extant literature tells us that politics is increasingly becoming
  nationalized, with more and more people taking their behavioral cues
  from the national political environment. However, we also know that
  there are still a plethora of ways in which local context shapes lived
  realities that should intersect with political behavior.
\end{itemize}

\begin{itemize}
\tightlist
\item
  That said, civic engagement is often conceptualized as coterminous
  with individual level participation, making it hard to disentangle
  underlying ideological considerations from community level variables
  that might shape access, information, and willingness to participate
  and the nature of that participation.
\end{itemize}
\end{frame}

\begin{frame}{Civic Opportunity Dataset}
\phantomsection\label{civic-opportunity-dataset}
\begin{figure}[H]

{\centering \pandocbounded{\includegraphics[keepaspectratio]{civicoppscreenshot.png}}

}

\caption{de Vries, Han, and Kim (\emph{Nature 2024})}

\end{figure}%
\end{frame}

\begin{frame}{Method}
\phantomsection\label{method}
\begin{itemize}
\tightlist
\item
  We merge the Civic Opportunity Dataset (henceforth COD) with the 2024
  wave of the Collaborative Multiracial Postelection Survey (henceforth
  CMPS) by connecting the county-level FIPS codes in the COD to the
  county level indicator in the CMPS. This allows us to examine the
  relationship between civil society and political behavior in the
  context of the 2024 election
\end{itemize}
\end{frame}

\begin{frame}{Research Questions}
\phantomsection\label{research-questions}
\begin{itemize}
\tightlist
\item
  In what ways did civic engagement at the local level shape the
  political decision-making of communities of color in the 2024
  election? We look at this in three ways:

  \begin{itemize}
  \tightlist
  \item
    Does the density of civil society in a given county predict the
    likelihood that a respondent in that county voted for Donald Trump
    in the 2024 election?
  \item
    Does the density of civil society in a given county predict partisan
    identification and strength for a given resident within that county?
  \item
    Does the density of civil society in a given county attenuate the
    likelihood that a respondent's racial resentment and/or
    anti-immigrant bias would predict their likelihood of voting for
    Donald Trump in the 2024 election?
  \end{itemize}
\end{itemize}
\end{frame}

\begin{frame}{Results}
\phantomsection\label{results}
\end{frame}




\end{document}
